\documentclass[../memoire.tex]{subfiles}
\begin{document}

\section*{Introduction}

\begin{comment}
Web security is fundamentals for nowadays enterprise applications and the security a user feels when registering or logging in to a website can leverage his experience and by extension his fidelity to the enterprise application / enterprise. To achieve secure authentication and authorization, many protocols have been developed. For example, SAML is a highly secured protocol for exchanging authentication and authorization data between parties. More over, OAuth is an authorization protocol that allows users to grant access to their resources on one site to another sit. And OpenID Connect is an authentication protocol built on top of OAuth 2.0 that allows clients to verify the identity of the end-user based on the authentication performed by an authorization server. 

Comparing this famous protocols is important to understand their strenghts and their weaknesses and to choose the right one for a specific background enterprise application. In this paper, we will compare SAML, OpenID Connect, and OAuth 2.0 for business applications by analyzing their features, security, and use cases. We will also discuss the advantages and disadvantages of each protocol and provide recommendations for choosing the right protocol for different business scenarios. Even in these protocols, there exist recommended practices to implement them in a secure way. We will also discuss these practices and how they can help to enhance the security of the protocols. 

The main questions we will try to answer in this paper are: What are the main differences between SAML, OpenID Connect, and OAuth 2.0? How do they work? Which protocol is best suited for different business scenarios? What are the recommended practices for implementing these protocols securely? 

Current developments privilege the use of OpenID Connect and OAuth 2.0 for modern web applications, but SAML is still widely used in enterprise applications. By comparing these protocols, we can help businesses make informed decisions about which protocol to use for their specific needs and how to implement it securely. 

This paper also discusses about their respective performance and scalability and how they fit in modern web application and Europian regulations such as GDPR. Finally, we will conclude by summarizing the main points of the comparison and providing recommendations for businesses looking to implement secure authentication and authorization protocols.
\end{comment}




Web security is a fundamental pillar of contemporary enterprise applications. The level of security a user experiences during registration or authentication can significantly enhance their user experience and, consequently, their loyalty to the platform or organization. To achieve secure authentication and authorization, various industry-standard protocols have been developed. For instance, Security Assertion Markup Language (SAML) is a mature, highly secure protocol designed for exchanging authentication and authorization data between distinct domains. In parallel, OAuth 2.0 serves as an authorization framework that allows users to grant limited access to their resources on one site to another without exposing credentials. Built on top of OAuth 2.0, OpenID Connect (OIDC) introduces a dedicated identity layer, enabling clients to verify the end-user's identity based on the authentication performed by an authorization server.

Conducting a comparative analysis of these widely adopted protocols is essential to understand their respective strengths and weaknesses, ultimately guiding the selection of the most appropriate framework for specific enterprise architectures. This paper evaluates SAML, OpenID Connect, and OAuth 2.0 within business environments by analyzing their technical features, security models, and primary use cases. Furthermore, we examine the core advantages and disadvantages of each protocol and formulate concrete recommendations tailored to diverse business scenarios. Because the security of these protocols heavily relies on proper deployment, this study also highlights established best practices and security profiles required to mitigate modern attack vectors.

To guide this research, this paper aims to address the following central questions: What are the architectural differences between SAML, OpenID Connect, and OAuth 2.0? How do their respective token and assertion handshakes function? Which protocol is best suited for specific corporate deployment scenarios? Finally, what are the recommended engineering practices to implement these protocols securely?

While current software engineering trends heavily favor the adoption of OpenID Connect and OAuth 2.0 for modern cloud-native web applications, SAML remains deeply entrenched in traditional enterprise infrastructures. By contrasting these protocols, this work provides businesses with the necessary insights to make informed architectural decisions based on their legacy constraints and future needs.

Additionally, this paper investigates the performance, data density, and scalability implications of each protocol, assessing how they align with modern web architectures and European regulatory frameworks such as the General Data Protection Regulation (GDPR). Finally, we conclude by synthesizing the core findings of this comparative study and outlining strategic recommendations for organizations aiming to deploy robust, compliant, and secure identity management solutions.

The different parts of this work are discussed in the next chapters.

\end{document}
